\chapter{Úvod}

Herní vývojáři se snaží vytvářet herní prostředí, které působí pro hráče co nejvíce uvěřitelně a interaktivně. 
Jedním z hlavních prostředků, jak tohoto cíle dosáhnout, jsou nehráčské postavy.

Nehráčská postava (z anglického \textit{Non-Player Character}, dále zkráceně NPC) představuje herní entitu, 
která není přímo ovládána hráčem, ale jejíž chování je řízeno umělou inteligencí. NPC pomáhají řešit problém 
„mrtvého světa“, tedy prostředí, které nereaguje na hráčovy akce a neposkytuje odpovídající zpětnou vazbu.

S rostoucí komplexitou herních světů se zvyšují také nároky na chování NPC. Jednoduché skriptované reakce často 
nestačí k vytvoření realistického dojmu a mohou působit repetitivně nebo nepřirozeně. Návrh komplexního chování 
NPC proto představuje významný problém, který je nutné řešit vhodnými metodami umělé inteligence.

Jedním z přístupů k implementaci chování NPC jsou behaviorální stromy, které umožňují rozdělit rozhodovací logiku do přehledné hierarchické struktury. 
Tento přístup je v současnosti běžně využíván v herním vývoji díky své  modularitě a schopnosti reagovat na změny prostředí.

Tato práce se zaměřuje na využití behaviorálních stromů v prostředí Unreal Engine 5 (UE5), což je moderní herní 
engine poskytující nástroje pro tvorbu herní umělé inteligence. Cílem práce je analyzovat principy behaviorálních 
stromů a demonstrovat jejich využití při implementaci chování NPC ve vzorové hře.

\medskip

Struktura práce je následující. Kapitola~2 se zabývá principy behaviorálních stromů a jejich vlastnostmi. 
V kapitole~3 jsou popsány nástroje Unreal Engine~5, které slouží k implementaci chování NPC. 
Kapitola~4 představuje návrh a architekturu vytvořeného systému NPC. 
V kapitole~5 je popsána samotná implementace jednotlivých komponent. 
Kapitola~6 se věnuje testování a evaluaci vytvořeného řešení.